% Colors for stuff
\definecolor{thmcolor}{RGB}{30, 144, 255}
\definecolor{lemcolor}{RGB}{50, 205, 50}
\definecolor{excolor}{RGB}{100, 80, 180}
\definecolor{defcolor}{RGB}{255, 140, 0}
\definecolor{propcolor}{RGB}{30,144,255}
\definecolor{rmkcolor}{RGB}{147, 112, 219}
\definecolor{corcolor}{RGB}{50, 205, 50}
\definecolor{notecolor}{RGB}{147, 112, 219}
\definecolor{quecolor}{RGB}{147, 112, 219}
\definecolor{exercolor}{RGB}{147, 112, 219}
\definecolor{conjcolor}{RGB}{30,144,255}

\newcounter{sharedcounter}

\tcbset{
    mybox/.style={
        arc=3pt,
        boxrule=0.5pt,
        titlerule=0.5pt,
        colback=white, 
        coltext=black,
        fonttitle=\bfseries,
        fontupper=\normalfont,
        top=4pt,
        bottom=4pt,
        left=6pt,
        right=6pt,
        toptitle=3pt,
        bottomtitle=3pt,
        boxsep=1pt,
        before skip=12pt,
        after skip=12pt,       
        breakable,
        label type=theorem,
        enhanced
    },
    plainbox/.style={
        boxrule=1pt,
        arc=0pt,
        colback=white,
        colframe=white,
        size=fbox,
        before skip=5pt,
        after skip=5pt
    },
    exbox/.style={
        sharp corners,
        % leftrule=-1pt,
        % rightrule=-1pt,
        colbacktitle=white,
        toprule at break=0pt,
        bottomrule at break=0pt,
        attach boxed title to top left={yshift=-0.5pt},
        boxed title style={
            sharp corners,
            boxrule=0.5pt,
            bottomrule=-0.5pt
        }
    }
}

% This defines a new box type with the same universal style, where the arguments reflect slight variations in style:
% 1. Environment name
% 2. Display name
% 3. General color theme
% 4. Cleverref Type
% 5. Cleverref Type Plural
% 6. Optional format changes
\NewDocumentCommand{\newboxtype}{m m m m m O{}}{
\NewTColorBox[use counter=sharedcounter, number within=section, crefname={#4}{#5}]{#1}{O{} O{}}
{mybox,
colframe=#3!50!black,
colbacktitle=#3!10,
coltitle=#3,
title=\ifblank{##1}{#2 \hfill \thetcbcounter}{\textit{##1} \hfill  \thetcbcounter},
##2,
#6
}
}

\newboxtype{thm}{Theorem}{thmcolor}{Theorem}{Theorems}
\newboxtype{defn}{Definition}{defcolor}{Definition}{Definitions}
\newboxtype{lem}{Lemma}{lemcolor}{Lemma}{Lemmas}
\newboxtype{cor}{Corollary}{corcolor}{Corollary}{Corollaries}
\newboxtype{prop}{Proposition}{propcolor}{Proposition}{Propositions}
\newboxtype{rmk}{Remark}{rmkcolor}{Remark}{Remarks}
% \newboxtype{ex}{Example}{excolor}{Example}{Examples}[exbox]
\newboxtype{que}{Question}{quecolor}{Question}{Questions}
\newboxtype{exer}{Exercise}{exercolor}{Exercise}{Exercises}
\newboxtype{conj}{Conjecture}{conjcolor}{Conjecture}{Conjectures}

\NewTColorBox[use counter=sharedcounter, number within=section, crefname={Example}{Examples}]{ex}{O{} O{}}
{mybox,
exbox,
colframe=excolor!50!black,
coltitle=excolor,
title=\hspace{-1.5pt}\ifblank{#1}{Example \thetcbcounter}{\textit{#1}},
#2
}

\NewTColorBox{note}{O{} O{}}
{mybox,
colframe=notecolor!50!black,
colbacktitle=notecolor!10,
coltitle=notecolor,
title=\ifblank{#1}{Note}{#1},
#2,
}

% An invisible box for centering floats inside other boxes
\newcommand{\figbox}[1]{
\begin{centering}
\tcbox[plainbox]{#1}
\end{centering}
}

\newtheoremstyle{boldExercise}
{6pt} % Space above
{6pt} % Space below
\bfseries % Body font
{} % Indent amount
{\color{Plum}\bfseries} % Head font
{:} % Punctuation after theorem head
{.5em} % Space after theorem head
{} % Theorem head spec 
\theoremstyle{boldExercise}
\newtheorem{exercise}{Exercise}